\input{../tex-inputs/latex-leseansicht-vorspann}

\section[Arthur und Olga Schnitzler an Gustav Schwarzkopf, 18. 6. 1909]{L04402 Arthur und Olga Schnitzler an Gustav Schwarzkopf, 18. 6. 1909}
\nopagebreak\mylabel{L04402v}
\rehead{ }\normalsize\beginnumbering\briefempfaengerindex{Schwarzkopf, Gustav@\textsc{Schwarzkopf, Gustav}!zzzSchnitzler, Olga@\emph{von Olga Schnitzler}!1909-06-182@{18. 6. 1909}|(be}\briefempfaengerindex{Schwarzkopf, Gustav@\textsc{Schwarzkopf, Gustav}!zzzSchnitzler, Arthur@\emph{von Arthur Schnitzler}!1909-06-182@{18. 6. 1909}|(be}
\toendnotes[C]{\smallbreak\pagebreak[2]}
\correspDesc{Versand  durch Arthur Schnitzler, Olga Schnitzler am 18. 6. 1909 in St. Gilgen
\newline{}Übermittlung  am 18. 6. 1909 in St. Gilgen
\newline{}Erhalt  durch Gustav Schwarzkopf im Zeitraum [19. 6. 1909 – 23. 6. 1909?] in Wien}\toendnotes[C]{\smallbreak}
\Standort{DLA, A:Schnitzler, HS.1985.1.1897.}
\physDesc{Bildpostkarte, 100 Zeichen
\newline{}Handschrift Arthur Schnitzler: Bleistift, deutsche Kurrent
\newline{}Handschrift Olga Schnitzler: schwarze Tinte
\newline{}Versand: Stempel: »\nobreak{}\oindex{St. Gilgen@\textbf{St. Gilgen}|pwk}St. Gilgen , 18. VI. 09\nobreak{}«.  }\pstart{}{\pb}Hrn \textsc{Gustav
                     Schwarzkopf}\pend{}\pstart{}\textsc{Wien I}\oindex{I., Innere Stadt@\textbf{I., Innere Stadt}|pw}\pend{}\pstart{}\textsc{Tiefer Graben 23}\oindex{Wien@\textbf{Wien}!I., Innere Stadt@\textbf{I., Innere Stadt}!Tiefer Graben 23@\textbf{Tiefer Graben 23}|pw}\pend{}{\bigskip}
\pstart
           \centering{}{\pb}\textcolor{gray}{\textbf{St. Gilgen (Salzkammergut)}}\oindex{St. Gilgen@\textbf{St. Gilgen}|pw}\pend
           \vspace{1em}
\pstart
           {\pb}18. 6. 07\pend
           \vspace{0.5em}
\pstart
           Herzliche Grüße.\pend
           \pstart Ihr \spacefill\mbox{Arthur}\pend{}\selectlanguage{ngerman}\vspace{1em}
\pstart
           \noindent{}{[}hs. O. Schnitzler:{]} Auf Wiedersehen!\pend
           \pstart \spacefill\mbox{Olga S.}\pend{}\selectlanguage{ngerman}\endnumbering\briefempfaengerindex{Schwarzkopf, Gustav@\textsc{Schwarzkopf, Gustav}!zzzSchnitzler, Olga@\emph{von Olga Schnitzler}!1909-06-182@{18. 6. 1909}|)be}\briefempfaengerindex{Schwarzkopf, Gustav@\textsc{Schwarzkopf, Gustav}!zzzSchnitzler, Arthur@\emph{von Arthur Schnitzler}!1909-06-182@{18. 6. 1909}|)be}\mylabel{L04402h}
\begin{anhang}
\end{anhang}\newcommand{\dateiname}{L04402}\newcommand{\titel}{Arthur und Olga Schnitzler an Gustav Schwarzkopf, 18. 6. 1909}\newcommand{\editorInnen}{Herausgegeben von Jahnke, SelmaMüller, Martin Anton}\input{../tex-inputs/latex-leseansicht-abspann}
